\documentclass[10pt,aspectratio=169]{beamer}

% ------------------------------------------------------------
% Presentacion 01 - Curso CGE y Microsimulacion Pensiones
% MACROPOL: Economia Aplicada y Politicas Publicas
% ------------------------------------------------------------

\usepackage[utf8]{inputenc}
\usepackage[T1]{fontenc}
\usepackage[spanish,es-nodecimaldot]{babel}
\usepackage{amsmath, amssymb, amsfonts}
\usepackage{booktabs}
\usepackage{array}
\usepackage{tikz}
\usepackage{xcolor}
\usepackage{graphicx}
\usepackage{hyperref}
\usepackage{ragged2e}
\usepackage{listings}

\usetikzlibrary{arrows.meta, positioning, shapes.geometric, calc}

% ------------------------------------------------------------
% Colores y estilo
% ------------------------------------------------------------
\definecolor{macropolblue}{RGB}{18,54,96}
\definecolor{macropolgreen}{RGB}{31,120,93}
\definecolor{macropolgray}{RGB}{90,90,90}
\definecolor{lightblue}{RGB}{232,240,249}
\definecolor{lightgreen}{RGB}{232,246,240}
\definecolor{lightgray}{RGB}{245,245,245}
\definecolor{accentorange}{RGB}{210,115,45}

\usetheme{Madrid}
\usecolortheme{default}
\setbeamercolor{structure}{fg=macropolblue}
\setbeamercolor{frametitle}{fg=white,bg=macropolblue}
\setbeamercolor{title}{fg=white,bg=macropolblue}
\setbeamercolor{block title}{fg=white,bg=macropolblue}
\setbeamercolor{block body}{fg=black,bg=lightblue}
\setbeamercolor{alerted text}{fg=accentorange}
\setbeamertemplate{navigation symbols}{}
\setbeamertemplate{itemize item}{\color{macropolgreen}$\blacktriangleright$}
\setbeamertemplate{itemize subitem}{\color{macropolgreen}$\circ$}
\setbeamertemplate{footline}[frame number]

\hypersetup{colorlinks=true, urlcolor=macropolblue, linkcolor=macropolblue}

\lstset{
  basicstyle=\ttfamily\scriptsize,
  backgroundcolor=\color{lightgray},
  frame=single,
  breaklines=true,
  columns=fullflexible
}

% ------------------------------------------------------------
% Macros simples
% ------------------------------------------------------------
\newcommand{\bigquestion}[1]{
\begin{center}
\Large\bfseries\color{macropolblue} #1
\end{center}
}

\newcommand{\source}[1]{\vspace{0.2cm}{\tiny\color{macropolgray} Fuente / referencia: #1}}

% ------------------------------------------------------------
% Datos de la presentacion
% ------------------------------------------------------------
\title[Curso CGE y Pensiones]{Modelos CGE y Microsimulaciones aplicados al Sistema de Pensiones de Capitalizacion Individual}
\subtitle{Presentacion 1: Roadmap del curso, motivacion internacional y arquitectura del modelo}
\author[MACROPOL]{Francisco Alberto Ramirez de Leon\\MACROPOL: Economia Aplicada y Politicas Publicas}
\date{2026}

\begin{document}

% ------------------------------------------------------------
\begin{frame}
  \titlepage
\end{frame}

% ------------------------------------------------------------
\begin{frame}{Proposito de la primera sesion}
\justifying
Esta primera sesion busca construir la intuicion general del curso antes de entrar en ecuaciones, codigo y calibracion.

\vspace{0.3cm}
\begin{block}{Idea central}
Una reforma previsional no es solo un cambio actuarial. Tambien afecta el mercado laboral, el ahorro, la inversion, el consumo, la distribucion del ingreso y la sostenibilidad fiscal.
\end{block}

\vspace{0.2cm}
Al final de la sesion, el participante debe entender por que necesitamos combinar:
\begin{itemize}
  \item modelos de Equilibrio General Computable, CGE;
  \item microsimulacion de hogares, trabajadores y cohortes;
  \item datos administrativos y encuestas de hogares;
  \item comunicacion de resultados para politica publica.
\end{itemize}
\end{frame}

% ------------------------------------------------------------
\begin{frame}{La pregunta que organiza todo el curso}
\bigquestion{Si reformamos el sistema de capitalizacion individual, ¿quien gana, quien pierde y que ocurre con la economia?}

\vspace{0.2cm}
Ejemplos de reformas:
\begin{itemize}
  \item aumentar la tasa de contribucion;
  \item modificar la edad de retiro;
  \item introducir una pension minima garantizada;
  \item reducir comisiones o cambiar reglas de inversion;
  \item fortalecer incentivos a la formalizacion.
\end{itemize}

\vspace{0.3cm}
\begin{alertblock}{Mensaje}
La calidad de la respuesta depende del modelo, de los datos y de la claridad con que se traduzca la institucion previsional en mecanismos economicos.
\end{alertblock}
\end{frame}

% ------------------------------------------------------------
\begin{frame}{Una reforma previsional tiene muchos canales}
\centering
\begin{tikzpicture}[node distance=1.25cm, every node/.style={font=\small}]
\tikzstyle{box}=[rectangle, rounded corners, draw=macropolblue, thick, fill=lightblue, align=center, minimum width=3.0cm, minimum height=0.75cm]
\tikzstyle{arrow}=[-Latex, thick, color=macropolblue]

\node[box, fill=lightgreen] (reform) {Reforma\\previsional};
\node[box, below left=of reform, xshift=-0.6cm] (labor) {Mercado\\laboral};
\node[box, below=of reform] (saving) {Ahorro y fondos\\de pensiones};
\node[box, below right=of reform, xshift=0.6cm] (fiscal) {Finanzas\\publicas};
\node[box, below=of saving, yshift=-0.2cm] (macro) {PIB, inversion,\\salarios y consumo};
\node[box, below=of macro] (micro) {Tasas de reemplazo,\\suficiencia y distribucion};

\draw[arrow] (reform) -- (labor);
\draw[arrow] (reform) -- (saving);
\draw[arrow] (reform) -- (fiscal);
\draw[arrow] (labor) -- (macro);
\draw[arrow] (saving) -- (macro);
\draw[arrow] (fiscal) -- (macro);
\draw[arrow] (macro) -- (micro);
\end{tikzpicture}

\vspace{0.2cm}
\small En equilibrio general, los efectos no ocurren aisladamente: se transmiten entre mercados, agentes e instituciones.
\end{frame}

% ------------------------------------------------------------
\begin{frame}{Herramientas disponibles para evaluar reformas}
\centering
\small
\begin{tabular}{p{3.0cm} p{4.3cm} p{4.8cm}}
\toprule
\textbf{Herramienta} & \textbf{Pregunta que responde} & \textbf{Limitacion principal} \\
\midrule
Modelo actuarial & ¿Cuanto pensiona una persona o cohorte? & No captura efectos macroeconomicos ni precios de equilibrio. \\
\addlinespace
Econometria & ¿Que relaciones historicas observamos en los datos? & No siempre permite simular reformas estructurales no observadas. \\
\addlinespace
Microsimulacion & ¿Quien gana y quien pierde dentro de la poblacion? & Requiere insumos macro: salarios, empleo, inflacion, retornos. \\
\addlinespace
CGE & ¿Como reaccionan mercados, precios, empleo, ahorro e inversion? & Agrega hogares y puede perder heterogeneidad distributiva. \\
\addlinespace
CGE + Micro & ¿Cual es el impacto macro, fiscal y distributivo integrado? & Requiere consistencia contable, datos y supuestos transparentes. \\
\bottomrule
\end{tabular}
\end{frame}

% ------------------------------------------------------------
\begin{frame}{Roadmap del curso}
\centering
\begin{tikzpicture}[every node/.style={font=\scriptsize}, node distance=0.85cm]
\tikzstyle{stage}=[rectangle, rounded corners, draw=macropolblue, thick, fill=lightblue, align=center, minimum width=2.45cm, minimum height=1.05cm]
\tikzstyle{arrow}=[-Latex, thick, color=macropolblue]

\node[stage, fill=lightgreen] (a) {A. Motivacion\\internacional};
\node[stage, right=of a] (b) {B. Herramientas\\R, Git, matrices};
\node[stage, right=of b] (c) {C. Microeconomia\\para CGE};
\node[stage, right=of c] (d) {D. SAM y\\CGE};
\node[stage, below=of d] (e) {E. Micro-\\simulacion};
\node[stage, left=of e] (f) {F. Integracion\\macro-micro};
\node[stage, left=of f, fill=lightgreen] (g) {G. Proyecto\\final};

\draw[arrow] (a) -- (b);
\draw[arrow] (b) -- (c);
\draw[arrow] (c) -- (d);
\draw[arrow] (d) -- (e);
\draw[arrow] (e) -- (f);
\draw[arrow] (f) -- (g);
\end{tikzpicture}

\vspace{0.35cm}
\begin{block}{Secuencia pedagogica}
Primero entendemos el problema de politica publica; luego las herramientas; despues la microeconomia; posteriormente la SAM, el CGE y la microsimulacion; finalmente integramos todo en una simulacion de reforma.
\end{block}
\end{frame}

% ------------------------------------------------------------
\begin{frame}{Estructura del curso por bloques}
\small
\begin{tabular}{p{2.0cm} p{4.1cm} p{6.0cm}}
\toprule
\textbf{Bloque} & \textbf{Tema} & \textbf{Resultado esperado} \\
\midrule
A & Evidencia internacional & Identificar preguntas de politica y resultados tipicos de reformas previsionales. \\
B & Herramientas & Preparar el entorno en R, organizar datos, matrices, funciones y repositorio. \\
C & Microeconomia & Entender consumidor, empresa, ahorro, oferta laboral y bienestar. \\
D & SAM y CGE & Construir una SAM simplificada y calibrar un CGE basico con pensiones. \\
E & Microsimulacion & Simular tasas de reemplazo, suficiencia y cobertura por grupos. \\
F & Integracion & Conectar resultados macro del CGE con resultados micro de hogares. \\
G & Proyecto final & Elaborar un informe tecnico de reforma previsional reproducible. \\
\bottomrule
\end{tabular}
\end{frame}

% ------------------------------------------------------------
\begin{frame}{Producto final del curso}
\begin{block}{Cada participante o equipo desarrollara una simulacion completa}
\begin{itemize}
  \item definicion de una reforma previsional;
  \item construccion o adaptacion de una SAM simplificada;
  \item calibracion de un modelo CGE basico;
  \item simulacion de impactos macroeconomicos;
  \item microsimulacion de tasas de reemplazo e indicadores distributivos;
  \item elaboracion de un policy brief tecnico.
\end{itemize}
\end{block}

\vspace{0.2cm}
\begin{alertblock}{Entrega esperada}
Un informe capaz de responder: \textbf{que cambia, por que cambia, cuanto cambia y para quienes cambia}.
\end{alertblock}
\end{frame}

% ------------------------------------------------------------
\begin{frame}{Por que mirar estudios internacionales}
\justifying
Antes de construir el modelo para el caso dominicano, conviene revisar que han respondido modelos similares en otros paises.

\vspace{0.25cm}
\begin{itemize}
  \item Permiten identificar canales economicos recurrentes.
  \item Ayudan a formular escenarios realistas.
  \item Muestran limites: informalidad, baja densidad de cotizacion, suficiencia y costo fiscal.
  \item Ayudan a explicar por que una reforma previsional es tambien una reforma macroeconomica.
\end{itemize}

\vspace{0.25cm}
\begin{block}{Instituciones que utilizan estas herramientas}
Banco Mundial, BID, OCDE, FMI, CEPAL, ministerios de hacienda, bancos centrales y superintendencias de pensiones.
\end{block}
\end{frame}

% ------------------------------------------------------------
\begin{frame}{Estudios internacionales: que problemas han respondido}
\small
\begin{tabular}{p{2.4cm} p{4.7cm} p{5.2cm}}
\toprule
\textbf{Caso} & \textbf{Pregunta de politica} & \textbf{Leccion para el curso} \\
\midrule
Chile & ¿Como afecta un sistema de cuentas individuales al ahorro, mercado laboral y suficiencia? & La capitalizacion individual puede fortalecer el ahorro, pero no resuelve por si sola baja densidad e informalidad. \\
\addlinespace
Suecia & ¿Como combinar sostenibilidad, edad de retiro y beneficios? & La edad de retiro y los ajustes actuariales son mecanismos centrales de sostenibilidad. \\
\addlinespace
Banco Mundial & ¿Como estructurar pilares de proteccion en vejez? & La suficiencia requiere combinar ahorro obligatorio, proteccion contra pobreza y ahorro voluntario. \\
\addlinespace
Modelos OLG & ¿Como afectan demografia y politica fiscal a generaciones distintas? & La reforma previsional es intertemporal e intergeneracional. \\
\bottomrule
\end{tabular}

\source{World Bank (1994); Auerbach y Kotlikoff; OECD, Pensions at a Glance; estudios de reforma previsional en Chile.}
\end{frame}

% ------------------------------------------------------------
\begin{frame}{Caso 1: Chile y cuentas individuales}
\begin{columns}[T]
\begin{column}{0.52\textwidth}
\textbf{Pregunta tipica}
\begin{itemize}
  \item ¿Que ocurre si aumenta la contribucion previsional?
  \item ¿Sube el ahorro agregado?
  \item ¿Mejoran las pensiones?
  \item ¿Hay efectos sobre empleo e informalidad?
\end{itemize}
\end{column}
\begin{column}{0.44\textwidth}
\begin{block}{Canales del modelo}
\begin{itemize}
  \item salario neto;
  \item costo laboral;
  \item ahorro obligatorio;
  \item inversion;
  \item acumulacion de fondos;
  \item pension futura.
\end{itemize}
\end{block}
\end{column}
\end{columns}

\vspace{0.2cm}
\source{Ejemplo motivacional basado en literatura sobre la reforma chilena y modelos dinamicos de equilibrio general aplicados a pensiones.}
\end{frame}

% ------------------------------------------------------------
\begin{frame}{Caso 2: Suecia y edad de retiro}
\begin{columns}[T]
\begin{column}{0.55\textwidth}
\textbf{Pregunta tipica}
\begin{itemize}
  \item ¿Que ocurre si aumenta la edad efectiva de retiro?
  \item ¿Como cambia la oferta laboral?
  \item ¿Como se ajusta la sostenibilidad del sistema?
\end{itemize}

\vspace{0.2cm}
\textbf{Resultado conceptual}
\begin{itemize}
  \item mas anos de cotizacion;
  \item menos anos de beneficio;
  \item mayor acumulacion individual;
  \item posible aumento del PIB por mayor empleo.
\end{itemize}
\end{column}
\begin{column}{0.4\textwidth}
\begin{alertblock}{Leccion}
La edad de retiro no es solo una variable actuarial. Tambien es una variable laboral y macroeconomica.
\end{alertblock}
\end{column}
\end{columns}

\source{OECD, Pensions at a Glance; notas pais sobre Suecia.}
\end{frame}

% ------------------------------------------------------------
\begin{frame}{Caso 3: pension minima o pilar solidario}
\begin{columns}[T]
\begin{column}{0.5\textwidth}
\textbf{Pregunta tipica}
\begin{itemize}
  \item ¿Que pasa si el Estado garantiza una pension minima?
  \item ¿Cuanto cuesta fiscalmente?
  \item ¿A quienes beneficia?
\end{itemize}
\end{column}
\begin{column}{0.46\textwidth}
\begin{block}{Aqui la microsimulacion es indispensable}
\begin{itemize}
  \item identifica hogares vulnerables;
  \item calcula brechas de suficiencia;
  \item mide incidencia por quintil;
  \item estima cambios en pobreza en vejez.
\end{itemize}
\end{block}
\end{column}
\end{columns}

\vspace{0.2cm}
\begin{alertblock}{Leccion}
El CGE muestra impactos agregados; la microsimulacion muestra distribucion, suficiencia y cobertura.
\end{alertblock}
\end{frame}

% ------------------------------------------------------------
\begin{frame}{Caso 4: informalidad y densidad de cotizacion}
\justifying
En America Latina, una reforma de capitalizacion individual puede estar bien disenada financieramente y aun asi producir pensiones insuficientes si la densidad de cotizacion es baja.

\vspace{0.25cm}
\begin{block}{Pregunta clave}
¿Que pasa si el trabajador no cotiza todos los meses de su vida laboral?
\end{block}

\vspace{0.2cm}
\begin{itemize}
  \item La tasa legal de contribucion no equivale a la tasa efectiva de acumulacion.
  \item La informalidad reduce fondos acumulados y cobertura.
  \item Los efectos de reforma dependen de comportamiento laboral y transiciones formal-informal.
\end{itemize}
\end{frame}

% ------------------------------------------------------------
\begin{frame}{Leccion general de los casos internacionales}
\centering
\begin{tikzpicture}[node distance=0.9cm, every node/.style={font=\small}]
\tikzstyle{box}=[rectangle, rounded corners, draw=macropolblue, thick, align=center, minimum width=4.0cm, minimum height=0.9cm]
\tikzstyle{arrow}=[-Latex, thick, color=macropolblue]

\node[box, fill=lightgreen] (policy) {Reforma previsional};
\node[box, below left=of policy, xshift=-1.4cm, fill=lightblue] (macro) {Efectos macro\\PIB, empleo, ahorro, inversion};
\node[box, below right=of policy, xshift=1.4cm, fill=lightblue] (micro) {Efectos micro\\cobertura, suficiencia, desigualdad};
\node[box, below=of policy, yshift=-2.0cm, fill=lightgreen] (eval) {Evaluacion integral\\CGE + microsimulacion};

\draw[arrow] (policy) -- (macro);
\draw[arrow] (policy) -- (micro);
\draw[arrow] (macro) -- (eval);
\draw[arrow] (micro) -- (eval);
\end{tikzpicture}

\vspace{0.3cm}
\small Ninguna reforma debe evaluarse solo por su impacto promedio: importa tambien la distribucion de los efectos.
\end{frame}

% ------------------------------------------------------------
\begin{frame}{Arquitectura macro-micro del curso}
\centering
\begin{tikzpicture}[node distance=0.85cm, every node/.style={font=\scriptsize}]
\tikzstyle{box}=[rectangle, rounded corners, draw=macropolblue, thick, fill=lightblue, align=center, minimum width=3.0cm, minimum height=0.8cm]
\tikzstyle{arrow}=[-Latex, thick, color=macropolblue]

\node[box, fill=lightgreen] (data) {Datos\\ENFT, ENIGH, administrativos};
\node[box, below left=of data, xshift=-0.8cm] (sam) {SAM\\matriz contable};
\node[box, below right=of data, xshift=0.8cm] (microdata) {Base micro\\trabajadores y hogares};
\node[box, below=of sam] (cge) {CGE\\precios y cantidades};
\node[box, below=of microdata] (microsim) {Microsimulacion\\cohortes y pensiones};
\node[box, below right=of cge, xshift=0.8cm, fill=lightgreen] (integration) {Integracion\\macro-micro};
\node[box, below=of integration, fill=lightgreen] (brief) {Policy brief\\decision publica};

\draw[arrow] (data) -- (sam);
\draw[arrow] (data) -- (microdata);
\draw[arrow] (sam) -- (cge);
\draw[arrow] (microdata) -- (microsim);
\draw[arrow] (cge) -- (integration);
\draw[arrow] (microsim) -- (integration);
\draw[arrow] (integration) -- (brief);
\end{tikzpicture}
\end{frame}

% ------------------------------------------------------------
\begin{frame}{Que es un modelo CGE}
\begin{block}{Definicion operativa}
Un modelo de Equilibrio General Computable es una representacion cuantitativa de una economia donde hogares, empresas, gobierno, sector externo y mercados interactuan simultaneamente hasta que oferta y demanda se igualan.
\end{block}

\vspace{0.2cm}
En el curso lo usaremos para simular:
\begin{itemize}
  \item cambios en contribuciones previsionales;
  \item acumulacion de fondos de pensiones;
  \item efectos sobre ahorro, inversion y capital;
  \item salarios, empleo y consumo;
  \item costo fiscal de medidas complementarias.
\end{itemize}
\end{frame}

% ------------------------------------------------------------
\begin{frame}{Que es una microsimulacion previsional}
\begin{block}{Definicion operativa}
Una microsimulacion previsional proyecta trayectorias individuales o de grupos de trabajadores para estimar aportes, fondos acumulados, pensiones y tasas de reemplazo bajo distintos escenarios.
\end{block}

\vspace{0.2cm}
Variables tipicas:
\begin{itemize}
  \item edad, sexo, educacion, quintil de ingreso;
  \item salario inicial y crecimiento salarial;
  \item densidad de cotizacion;
  \item tasa de contribucion;
  \item retorno financiero;
  \item edad de retiro;
  \item expectativa de vida o factor actuarial.
\end{itemize}
\end{frame}

% ------------------------------------------------------------
\begin{frame}{Por que integrar CGE y microsimulacion}
\small
\begin{tabular}{p{3.5cm} p{4.4cm} p{4.4cm}}
\toprule
 & \textbf{CGE} & \textbf{Microsimulacion} \\
\midrule
Unidad de analisis & Economia completa & Persona, hogar o cohorte \\
Precios & Endogenos & Usualmente exogenos \\
Heterogeneidad & Limitada & Alta \\
Fortaleza & Consistencia macro y de mercados & Distribucion e incidencia \\
Debilidad & Puede ocultar desigualdad & Puede ignorar retroalimentacion macro \\
Integracion & Salarios, empleo, retornos e inflacion & Ahorro, beneficios, cobertura y suficiencia \\
\bottomrule
\end{tabular}

\vspace{0.3cm}
\begin{alertblock}{Idea clave}
El CGE responde \textit{que ocurre con la economia}; la microsimulacion responde \textit{a quienes les ocurre}.
\end{alertblock}
\end{frame}

% ------------------------------------------------------------
\begin{frame}{Sistema de capitalizacion individual dentro del CGE}
\centering
\begin{tikzpicture}[node distance=0.8cm, every node/.style={font=\small}]
\tikzstyle{box}=[rectangle, rounded corners, draw=macropolblue, thick, fill=lightblue, align=center, minimum width=3.4cm, minimum height=0.8cm]
\tikzstyle{arrow}=[-Latex, thick, color=macropolblue]

\node[box, fill=lightgreen] (worker) {Trabajador formal};
\node[box, right=of worker] (contrib) {Contribucion\\previsional};
\node[box, right=of contrib] (fund) {Fondo AFP};
\node[box, below=of fund] (capital) {Mercado de capitales\\e inversion};
\node[box, left=of capital] (production) {Empresas\\produccion};
\node[box, left=of production] (income) {Salarios y\\rentas};

\draw[arrow] (worker) -- (contrib);
\draw[arrow] (contrib) -- (fund);
\draw[arrow] (fund) -- (capital);
\draw[arrow] (capital) -- (production);
\draw[arrow] (production) -- (income);
\draw[arrow] (income) -- (worker);
\end{tikzpicture}

\vspace{0.3cm}
\small La AFP no es solo un receptor contable de aportes: es un canal entre trabajo, ahorro, mercado financiero, inversion y pensiones futuras.
\end{frame}

% ------------------------------------------------------------
\begin{frame}{Ecuacion conceptual de acumulacion previsional}
\bigquestion{$PF_{i,t+1} = (1+r_t)PF_{i,t} + c_t \cdot w_{i,t} \cdot D_{i,t}$}

\vspace{0.2cm}
Donde:
\begin{itemize}
  \item $PF_{i,t}$: fondo previsional del individuo $i$;
  \item $r_t$: retorno financiero del fondo;
  \item $c_t$: tasa de contribucion;
  \item $w_{i,t}$: salario cotizable;
  \item $D_{i,t}$: indicador o densidad de cotizacion efectiva.
\end{itemize}

\vspace{0.2cm}
\begin{alertblock}{Conexion macro-micro}
El CGE puede generar $w_t$, $r_t$ y empleo; la microsimulacion aplica esos cambios a individuos o cohortes.
\end{alertblock}
\end{frame}

% ------------------------------------------------------------
\begin{frame}{El bloque de herramientas del curso}
\begin{columns}[T]
\begin{column}{0.48\textwidth}
\textbf{Herramientas computacionales}
\begin{itemize}
  \item R y RStudio;
  \item estructura de proyectos;
  \item Git y GitHub;
  \item lectura, limpieza y transformacion de datos;
  \item visualizacion de resultados.
\end{itemize}
\end{column}
\begin{column}{0.48\textwidth}
\textbf{Herramientas matematicas}
\begin{itemize}
  \item algebra matricial;
  \item sistemas de ecuaciones;
  \item optimizacion numerica;
  \item calibracion;
  \item analisis de escenarios.
\end{itemize}
\end{column}
\end{columns}

\vspace{0.25cm}
\begin{block}{Objetivo}
Que el participante no solo entienda el modelo, sino que pueda reconstruirlo, auditarlo y adaptarlo.
\end{block}
\end{frame}

% ------------------------------------------------------------
\begin{frame}[fragile]{Toolkit computacional MACROPOL}
\begin{lstlisting}
Curso_CGE_Pensiones/
|
|-- data/
|   |-- raw/
|   |-- processed/
|
|-- sam/
|   |-- 01_construccion_sam.R
|
|-- cge/
|   |-- 02_calibracion_cge.R
|   |-- 03_simulacion_reforma.R
|
|-- microsim/
|   |-- 04_tasas_reemplazo.R
|
|-- integration/
|   |-- 05_integracion_macro_micro.R
|
|-- outputs/
|-- policy_brief/
|-- README.md
\end{lstlisting}

\vspace{0.1cm}
\small Esta estructura convierte el curso en un laboratorio reproducible de evaluacion de reformas previsionales.
\end{frame}

% ------------------------------------------------------------
\begin{frame}{Microeconomia necesaria antes del CGE}
\justifying
Antes de calibrar el modelo, haremos un repaso aplicado de microeconomia. No sera un curso teorico completo, sino un puente entre intuicion economica y modelizacion computacional.

\vspace{0.25cm}
\begin{itemize}
  \item consumidor: consumo, ocio, ahorro y bienestar;
  \item empresa: funcion de produccion y demanda de factores;
  \item mercado laboral: salario neto, costo laboral y formalidad;
  \item ahorro intertemporal: consumo presente vs. consumo futuro;
  \item equilibrio: precios que ajustan oferta y demanda;
  \item bienestar: eficiencia, suficiencia y equidad.
\end{itemize}
\end{frame}

% ------------------------------------------------------------
\begin{frame}{La SAM: la fotografia contable de la economia}
\begin{block}{Definicion}
La Matriz de Contabilidad Social, SAM, organiza los flujos de ingresos y gastos entre sectores productivos, factores, hogares, gobierno, AFP, inversion y resto del mundo.
\end{block}

\vspace{0.2cm}
En el curso construiremos una SAM simplificada con cuentas para:
\begin{itemize}
  \item sectores productivos;
  \item trabajo y capital;
  \item hogares por grupo;
  \item gobierno;
  \item sistema previsional / AFP;
  \item ahorro-inversion;
  \item resto del mundo.
\end{itemize}
\end{frame}

% ------------------------------------------------------------
\begin{frame}{Reformas que simularemos}
\begin{columns}[T]
\begin{column}{0.32\textwidth}
\begin{block}{Parametricas}
\begin{itemize}
  \item contribucion;
  \item edad de retiro;
  \item comision;
  \item tasa de reemplazo objetivo.
\end{itemize}
\end{block}
\end{column}
\begin{column}{0.32\textwidth}
\begin{block}{Estructurales}
\begin{itemize}
  \item pilar solidario;
  \item pension minima;
  \item subsidio a cotizacion;
  \item sistema mixto.
\end{itemize}
\end{block}
\end{column}
\begin{column}{0.32\textwidth}
\begin{block}{Escenarios}
\begin{itemize}
  \item menor crecimiento;
  \item menor retorno;
  \item mayor informalidad;
  \item envejecimiento.
\end{itemize}
\end{block}
\end{column}
\end{columns}
\end{frame}

% ------------------------------------------------------------
\begin{frame}{Indicadores de salida}
\begin{columns}[T]
\begin{column}{0.48\textwidth}
\textbf{Macroeconomicos}
\begin{itemize}
  \item PIB;
  \item consumo;
  \item inversion;
  \item ahorro nacional;
  \item empleo formal e informal;
  \item salarios;
  \item acumulacion de capital.
\end{itemize}
\end{column}
\begin{column}{0.48\textwidth}
\textbf{Previsionales y distributivos}
\begin{itemize}
  \item fondo acumulado;
  \item pension esperada;
  \item tasa de reemplazo;
  \item cobertura;
  \item suficiencia;
  \item pobreza en vejez;
  \item incidencia por quintil y genero.
\end{itemize}
\end{column}
\end{columns}
\end{frame}

% ------------------------------------------------------------
\begin{frame}{Cronograma sugerido de 15 sesiones}
\scriptsize
\begin{tabular}{p{1.1cm} p{4.2cm} p{6.7cm}}
\toprule
\textbf{Sesion} & \textbf{Bloque} & \textbf{Contenido} \\
\midrule
1 & Motivacion y roadmap & Evidencia internacional, preguntas de politica, arquitectura CGE-micro. \\
2 & Herramientas I & R, estructura del proyecto, limpieza de datos, visualizacion. \\
3 & Herramientas II & Algebra matricial, sistemas, optimizacion y Git. \\
4 & Microeconomia I & Consumidor, ahorro intertemporal y oferta laboral. \\
5 & Microeconomia II & Empresa, produccion, equilibrio y bienestar. \\
6 & SAM I & Cuentas, flujos, balance contable y estructura previsional. \\
7 & SAM II & Construccion de SAM simplificada en R. \\
8 & CGE I & Ecuaciones, calibracion y cierre del modelo. \\
9 & CGE II & Simulacion de aumento de contribuciones. \\
10 & CGE III & Escenarios macrofinancieros y sensibilidad. \\
11 & Microsimulacion I & Trayectorias laborales, salarios y densidad de cotizacion. \\
12 & Microsimulacion II & Tasas de reemplazo y suficiencia por grupos. \\
13 & Integracion & Top-down, bottom-up e integracion hibrida. \\
14 & Comunicacion & Graficos, policy brief y presentacion de resultados. \\
15 & Proyecto final & Presentacion y discusion de simulaciones. \\
\bottomrule
\end{tabular}
\end{frame}

% ------------------------------------------------------------
\begin{frame}{Como leer resultados de modelos}
\begin{block}{Tres preguntas de interpretacion}
\begin{enumerate}
  \item \textbf{Mecanismo:} ¿por que cambia la variable?
  \item \textbf{Magnitud:} ¿cuanto cambia respecto al escenario base?
  \item \textbf{Incidencia:} ¿quien se beneficia y quien enfrenta costos?
\end{enumerate}
\end{block}

\vspace{0.25cm}
\begin{alertblock}{Advertencia metodologica}
Un resultado de CGE no es una prediccion mecanica. Es un experimento condicionado por supuestos, datos, elasticidades, cierre macroeconomico y escenario de referencia.
\end{alertblock}
\end{frame}

% ------------------------------------------------------------
\begin{frame}{Lecturas y referencias iniciales}
\small
\begin{itemize}
  \item World Bank (1994). \textit{Averting the Old Age Crisis: Policies to Protect the Old and Promote Growth}.
  \item Auerbach, A. y Kotlikoff, L. \textit{Dynamic Fiscal Policy}.
  \item OECD. \textit{Pensions at a Glance} y notas pais sobre reformas previsionales.
  \item Estudios aplicados de modelos dinamicos de equilibrio general para reformas previsionales en Chile.
  \item Documentos nacionales sobre el sistema de capitalizacion individual y datos laborales de encuestas de hogares.
\end{itemize}

\vspace{0.3cm}
\begin{block}{Uso en el curso}
Estas lecturas no se usaran como teoria abstracta, sino como ejemplos de preguntas, escenarios, mecanismos y formas de presentar resultados.
\end{block}
\end{frame}

% ------------------------------------------------------------
\begin{frame}{Mensaje de cierre}
\bigquestion{El objetivo no es aprender software.}

\bigquestion{El objetivo es transformar una reforma previsional en un experimento cuantitativo reproducible.}

\vspace{0.3cm}
\begin{center}
\large\color{macropolgreen}\textbf{Datos $\rightarrow$ modelo $\rightarrow$ resultados $\rightarrow$ decision publica}
\end{center}
\end{frame}

% ------------------------------------------------------------
\begin{frame}
\centering
\Huge\color{macropolblue}\textbf{Gracias}

\vspace{0.5cm}
\Large Preguntas y comentarios

\vspace{0.8cm}
\normalsize MACROPOL: Economia Aplicada y Politicas Publicas
\end{frame}

\end{document}
